\begin{lemma}\label{lem:decic54-gonality}
The curve $X_1(54)$ has rational gonality at least $12$.
\end{lemma}
\begin{proof}
Put $X=X_1(54)$ and let $Y$ be the quotient by the order-$3$
projective diamond $\langle17\rangle$.
Then $g(X)=52$ and $g(Y)=10$.
For completeness, the congruence subgroup of $Y$ has index $324$,
no elliptic points, and $36$ cusps. Thus its genus is
$1+324/12-36/2=10$.
The cusp count can also be obtained directly: $X$ has $60$ geometric
cusps, of which $24$ are fixed by $\langle17\rangle$, so the quotient
has $(60+2\cdot24)/3=36$ cusps.
The exact primitive-coset calculation is recorded in
\path{parent_gonality_controls.json}.

The analytic bound \cite[(2.1)]{NV} gives
\[
 \operatorname{gon}_{\mathbb C}(X)
 \geq\frac{325\cdot972}{32768}>9.
\]
Suppose that $f:X\longrightarrow\mathbb P^1$ has degree
$n\in\{10,11\}$. Since $[\mathbb Q(X):\mathbb Q(Y)]=3$ is prime,
either $f$ factors through $Y$, or the two maps generate
$\mathbb Q(X)$. The first alternative would imply $3\mid n$.
In the second, Castelnuovo--Severi gives
\[
 52=g(X)\leq 3g(Y)+(3-1)(n-1)\leq50,
\]
a contradiction. This excludes both remaining degrees.
\end{proof}

\begin{lemma}\label{lem:decic54-good-ordinary2}
If $E/K$ has a point of order $54$ and $[K:\mathbb Q]=10$,
then $E$ has good ordinary reduction at some prime above $2$.
At every good prime above $2$ the ramification index is $1$ and
the residue degree is between $6$ and $10$.
\end{lemma}
\begin{proof}
The rational Jacobian $J_1(54)(\mathbb Q)$ has rank zero
\cite[Theorem~3.1(2)]{DEHMZ}. Its rational cuspidal subgroup has
invariants $[3,3,3,9,9,1102437,1529080119]$ and index dividing $3$
\cite[Theorem~4.13 and Table~2]{DEHMZ}; hence the whole group has
odd order.
At $2$, the regular semistable model consists of two copies of
$X_1(27)_{\mathbb F_2}$, with completed strict local equations
$xy-2$. The auxiliary $27$-structure removes automorphisms.
The cuspidal boundary is smooth and finite etale, with the same
$60$ geometric points in the generic and special fibres.
Its global cusp signature is $1^{18}2^6 6^2 9^2$.
Since $2$ generates $(\mathbb Z/27\mathbb Z)^\times$, all its
closed special cusp orbits lift to rational cuspidal divisors of
the same degrees.

Suppose that every prime above $2$ were bad. The point of order
$27$ excludes additive reduction. The degree-$10$ trace divisor
$D$ on $X_1(54)$ would therefore specialize to the smooth cuspidal
boundary. Lift it to a rational cusp divisor $C$ of degree $10$.
The Abel map on the smooth locus extends to the Neron model,
and $[D-C]$ reduces to zero in its entire special fibre.
That kernel is pro-$2$, whereas $J_1(54)(\mathbb Q)$ has odd order.
Thus $D\sim C$, contradicting Lemma~\ref{lem:decic54-gonality}.

At a good prime the point of order $27$ survives. The exact
noncuspidal closed-place counts on $X_1(27)_{\mathbb F_2}$ vanish
through degree $5$, so its residue degree $f$ is at least $6$.
From $ef\leq10$ we obtain $e=1$.
If the reduction were supersingular, a rational point of order
$2$ would lie in the formal group. In an integral Weierstrass
model $a_1$ is even, and
\[
 [2](T)=2T-a_1T^2+O(T^3).
\]
For a nonzero $T$ over an unramified $2$-adic field,
$v_2(T)\geq1$. The first displayed term has strictly smaller
valuation than every other term, so $[2](T)\ne0$.
This contradiction proves ordinarity.
\end{proof}


\begin{proposition}\label{prop:decic54}
There is no elliptic curve over a number field of degree $10$ with a point
of order $54$.
\end{proposition}
\begin{proof}
Put $X=X_1(54)$ and $Z=X_1(27)$. The maps
\[
 \alpha(E,P_2,Q)=(E,Q),\qquad
 \beta(E,P_2,Q)=(E/\langle P_2\rangle,\overline Q)
\]
are defined over $\mathbf Q$. Here $P_2$ has order $2$ and $Q$ has
order $27$. Let $D$ be the degree-$10$ trace divisor supplied by a
hypothetical point of order $54$, and write
$D_0=\alpha_*D$, $D_1=\beta_*D$. Proper pushforward is used throughout,
including all residue-degree multiplicities.

The rational gonality of $X$ is at least $12$, and $J_X(\mathbf Q)$
is finite; see Lemma~\ref{lem:decic54-gonality} and
\cite[Theorem~3.1(2)]{DEHMZ}. At any odd good prime, an entirely
cuspidal reduction of $D$ is impossible. Indeed, lift that reduction
to a cuspidal divisor over the local field. Its difference from $D$
is torsion by Manin--Drinfeld and the rank-zero assertion, and has
zero reduction. Torsion injectivity at an odd prime makes the
difference principal. The two effective divisors are distinct,
so base change would give $h^0(D)\geq2$ over $\mathbf Q$,
contrary to the gonality bound.

By Lemma~\ref{lem:decic54-good-ordinary2}, there is a good ordinary
prime above $2$, unramified and of residue degree $f\in\{6,\ldots,10\}$.
It is the unique good prime above $2$, since two such residue degrees
would exceed $10$. Replacing the elliptic curve by its quotient by
$P_2$, if necessary, we may arrange that the chosen degree-$2$
isogeny has connected kernel there. The replacement again has a
rational point of order $2$ and a point of order $27$.
On the ordinary component the two degeneracy maps then differ by
Frobenius, so their trace reductions have the same noncuspidal
closed point $N$ of degree $f$, with multiplicity one. On the two
cuspidal components they agree or differ by $\delta=\langle2\rangle$.
Consequently, for effective rational cuspidal divisors $C_a,C_e$ on $Z$,
\[
 \overline D_0=N+\overline C_a+\overline C_e,\qquad
 \overline D_1=N+\overline C_a+\delta\overline C_e,
 \qquad \deg(C_a+C_e)=10-f.                         \tag{54.1}
\]
All cusps of degrees at most $4$ at $2$ are globally intact, so the
cuspidal divisors in this formula lift over $\mathbf Q$ with the
same degrees.

The common cusp identification can be made integrally. In Tate normal
form write
\[
 B=16b^2+(1-20c-8c^2)b+c(c-1)^3.
\]
The functions $b$ and $b/B$ are units on the open integral modular curve:
the discriminant is $b^3B$, up to sign. Their cusp valuation vectors
therefore commute with good reduction. There are $14$ global cusp orbits,
of degrees $1^9,2^3,6,9$. Both $2$ and $5$ generate
$(\mathbf Z/27\mathbf Z)^\times$, so these orbits are intact at both
primes. Their $14$ valuation vectors are distinct and agree in the two
characteristics. In particular, the $12$ cusps of degrees at most $2$
are identified across the two reductions, as is their diamond action.

The group
\[
 J_Z(\mathbf Q)\simeq C_3^2\times C_{52497}
\]
is generated by rational cuspidal divisors
\cite[Theorem~4.13 and Table~2]{DEHMZ}. Its order is odd, and reduction
at $2$ identifies it with $J_Z(\mathbf F_2)$. In particular,
\[
 D_1-D_0\sim\delta C_e-C_e.                         \tag{54.2}
\]
If
\[
 h^0\bigl(N+\overline C_a+\overline C_e+
                     \delta\overline C_e\bigr)=1,                 \tag{54.3}
\]
then semicontinuity, applied to the equivalent effective divisors
$D_1+C_e$ and $D_0+\delta C_e$, makes these divisors equal.
Their noncuspidal parts give $D_1=D_0$.
This is impossible. Equality would make the marked degree-$2$
isogenies permute the at most $10$ geometric points in the support
of $D_0$. A cycle of length $r\leq10$ would give an endomorphism of
degree $2^r$ carrying a marked point of order $27$ to itself up to sign.
For a non-CM elliptic curve this endomorphism is multiplication by
$\pm2^{r/2}$, so $r$ is even and $2^{r/2}\equiv\pm1\pmod{27}$.
The least positive exponent with this property is $9$, a contradiction.
The CM case is excluded by the minimum-degree table of \cite{CGPS}:
the least degree of a CM point on $X_1(54)$ is $18$.
The exact table row is retained in \path{inputs/CM54_minimum_check.json}.

We compute (54.3) for every possibility in (54.1). The noncuspidal
closed points of $Z/\mathbf F_2$ vanish through degree $5$.
After simultaneous diamonds there are $25$ ordinary noncuspidal representatives
in degrees $6$ through $10$. All colored cusp complements of degree
$10-f$ are included. Of the resulting $8825$ cases, $4721$ satisfy
(54.3). The remaining $4104$ cases give $4103$ ordered pairs of rational
Jacobian classes, with $800$ distinct first classes. These lists, the
chosen rational cusp $A$, and all cusp-class coordinates are recorded in
\path{dyadic_rigidity_v2.m}. Classes are normalized by subtracting $10A$.

We test these $800$ classes at $5$. The characteristic-zero cusp
identifications are retained in both characteristics. For each rational
class $\xi$ we construct a divisor in the class $10A+\xi$ over
$\mathbf F_5$ and compute its complete Riemann--Roch space. Exactly
$724$ of the $800$ spaces are nonzero. Enumerating every projective
section gives $1067$ effective divisors with a noncuspidal part and
$217$ entirely cuspidal divisors. The latter are impossible by the
first paragraph, since $\alpha^{-1}$ of the cusp locus is the cusp
locus. There are $192$ distinct noncuspidal closed points in the
remaining divisors.

For each of these points we factor the exact $2$-division cubic of
the represented elliptic curve. An irreducible factor of degree $e$
corresponds to a Frobenius orbit of $2$-kernels of length $e$.
If the point occurs with multiplicity $m$, all distributions
$\sum e_i n_i=m$ are included. The quotient isogenies retain the
marked point of order $27$ and the proper pushforward norm of each
image. There are $361$ kernel orbits and $334$ distinct marked target
images. Only $147$ complete lifts of the noncuspidal parts are possible:
$59$ have degree $10$, and $88$ have degree $9$ with a cusp complement
of degree $1$.

For every lift we compute the class of its second degeneracy image.
We then allow every effective cuspidal complement of the required
degree on $Z/\mathbf F_5$, without imposing a lift condition on that
complement. None of the $147$ lifts gives one of the $4103$ ordered
rational class pairs allowed by (54.1). The exact final test is recorded
in \path{member_all_classes.m}. Thus no reduction of $D$ is possible,
which proves the proposition.
\end{proof}
